\newcommand{\C}{\mathbb{C}}
\newcommand{\R}{\mathbb{R}}
\newcommand{\Z}{\mathbb{Z}}
\newcommand{\Tr}{\operatorname{Tr}}
\newcommand{\Impart}{\operatorname{Im}}
\newcommand{\AF}{\operatorname{AF}}
\newcommand{\AX}{\operatorname{AX}}
\newcommand{\GL}{\operatorname{GL}}

\examyear{2012}
% \examera{平成24年}
% \examfield{代数}
% \examgroup{B}
% \examnumber{4}
% \examtitle{平成24年 B 第4問}
% \examtag{非可換代数}
% \examtag{表現論}
% \examtag{既約表現}

\(A\) を \(2\) 元 \(e,f\) で生成される単位元 \(1\) を持つ複素数体 \(\C\) 上の結合的自由代数とする。
\(I\) を \(e^2-e,\ f-ef-fe\) で生成される \(A\) の両側イデアルとし，商代数 \(B=A/I\) を考える。

\begin{enumerate}
  \item \(a_1=ef,\ a_2=(1-e)f\) とするとき，次の等式を示せ。
  \[
    a_1e=0,\qquad a_2(1-e)=0.
  \]
  \item \(B\) の \(\C\) 上の \(1\) 次元表現を全て求めよ。
  \item \(B\) の \(\C\) 上の有限次元既約表現を全て求めよ。ただし，\(B\) の表現 \(V\) が既約であるとは，
  \(V\ne\{0\}\) かつ \(V\) が自分自身と \(\{0\}\) 以外に部分表現を持たないことをいう。
\end{enumerate}

%